\documentclass[10pt,a4paper]{amsart}
\usepackage{amsmath,amssymb,amsthm,graphicx}
% For commutative diagrams we use Xy-pic
\usepackage[all]{xy}

% We use 2cell for 2-commutative diagrams.
\xyoption{2cell}
% Theorem environments.
%
\theoremstyle{plain}
\newtheorem{theorem}[subsection]{Theorem}
\newtheorem{proposition}[subsection]{Proposition}
\newtheorem{lemma}[subsection]{Lemma}

\theoremstyle{definition}
\newtheorem{definition}[subsection]{Definition}
\newtheorem{example}[subsection]{Example}
\newtheorem{exercise}[subsection]{Exercise}
\newtheorem{situation}[subsection]{Situation}

\theoremstyle{remark}
\newtheorem{remark}[subsection]{Remark}
\newtheorem{remarks}[subsection]{Remarks}

\numberwithin{equation}{subsection}

% Macros
%
\def\lim{\mathop{\mathrm{lim}}\nolimits}
\def\colim{\mathop{\mathrm{colim}}\nolimits}
\def\Spec{\mathop{\mathrm{Spec}}}
\def\Hom{\mathop{\mathrm{Hom}}\nolimits}
\def\Ext{\mathop{\mathrm{Ext}}\nolimits}
\def\SheafHom{\mathop{\mathcal{H}\!\mathit{om}}\nolimits}
\def\SheafExt{\mathop{\mathcal{E}\!\mathit{xt}}\nolimits}
\def\Sch{\mathit{Sch}}
\def\Mor{\mathop{\mathrm{Mor}}\nolimits}
\def\Ob{\mathop{\mathrm{Ob}}\nolimits}
\def\Sh{\mathop{\mathit{Sh}}\nolimits}
\def\NL{\mathop{N\!L}\nolimits}
\def\CH{\mathop{\mathrm{CH}}\nolimits}
\def\proetale{{pro\text{-}\acute{e}tale}}
\def\etale{{\acute{e}tale}}
\def\QCoh{\mathit{QCoh}}
\def\Ker{\mathop{\mathrm{Ker}}}
\def\Im{\mathop{\mathrm{Im}}}
\def\Coker{\mathop{\mathrm{Coker}}}
\def\Coim{\mathop{\mathrm{Coim}}}

% Boxtimes
%
\DeclareMathSymbol{\boxtimes}{\mathbin}{AMSa}{"02}

%
% Macros for moduli stacks/spaces
%
\def\QCohstack{\mathcal{QC}\!\mathit{oh}}
\def\Cohstack{\mathcal{C}\!\mathit{oh}}
\def\Spacesstack{\mathcal{S}\!\mathit{paces}}
\def\Quotfunctor{\mathrm{Quot}}
\def\Hilbfunctor{\mathrm{Hilb}}
\def\Curvesstack{\mathcal{C}\!\mathit{urves}}
\def\Polarizedstack{\mathcal{P}\!\mathit{olarized}}
\def\Complexesstack{\mathcal{C}\!\mathit{omplexes}}
% \Pic is the operator that assigns to X its picard group, usage \Pic(X)
% \Picardstack_{X/B} denotes the Picard stack of X over B
% \Picardfunctor_{X/B} denotes the Picard functor of X over B
\def\Pic{\mathop{\mathrm{Pic}}\nolimits}
\def\Picardstack{\mathcal{P}\!\mathit{ic}}
\def\Picardfunctor{\mathrm{Pic}}
\def\Deformationcategory{\mathcal{D}\!\mathit{ef}}
\newenvironment{reference}{}{}
\newenvironment{slogan}{}{}
\newenvironment{history}{}{}
\usepackage{cleveref}
\crefname{originallemma}{original lemma}{original lemmas}
\Crefname{originallemma}{Original lemma}{Original lemmas}
\makeatletter
\newlabel{lemma-genus-zero}{{0C8S}{source}{}{sourcepointer-1}{}}
\newlabel{lemma-genus-zero@cref}{{[originallemma][0][]0C8S}{[1][0][]source}{}{}{}}
\makeatother
\crefname{originallemma}{original lemma}{original lemmas}
\Crefname{originallemma}{Original lemma}{Original lemmas}
\makeatletter
\newlabel{item-one}{{0C8U}{source}{}{sourcepointer-2}{}}
\newlabel{item-one@cref}{{[originallemma][0][]0C8U}{[1][0][]source}{}{}{}}
\makeatother
\crefname{originallemma}{original lemma}{original lemmas}
\Crefname{originallemma}{Original lemma}{Original lemmas}
\makeatletter
\newlabel{lemma-genus-one}{{0C8T}{source}{}{sourcepointer-3}{}}
\newlabel{lemma-genus-one@cref}{{[originallemma][0][]0C8T}{[1][0][]source}{}{}{}}
\makeatother
\crefname{originallemma}{original lemma}{original lemmas}
\Crefname{originallemma}{Original lemma}{Original lemmas}
\makeatletter
\newlabel{lemma-recurring-symmetric-real}{{0C5X}{source}{}{sourcepointer-4}{}}
\newlabel{lemma-recurring-symmetric-real@cref}{{[originallemma][0][]0C5X}{[1][0][]source}{}{}{}}
\makeatother
\title{Genus dominates graph genus for a minimal numerical type}
\author{The Stacks Project Authors}
\date{stacks-project:models.tex:local-snapshot-20261011; GNU FDL 1.2 or later}
\begin{document}
\maketitle
\noindent\textit{Editorial note.} Original author theorem and proof are retained. Only this target is counted; local supporting statements are not separate new entries.\par\medskip
\noindent\textit{Allowed standard prerequisites (editorial).}\begin{itemize}
\item Elementary graph homology, weighted symmetric intersection matrices and the explicit definitions of numerical type, minimality and genus retained in the source.. These are the named standard inputs; the selected main construction and its local nonroutine support remain in full.
\end{itemize}\medskip
\ifdefined\hypertarget\hypertarget{sourcepointer-1}{}\fi
\noindent\textit{Source reference.} Original lemma 0C8S: focused-seeds/AGP/sources/stacks-snapshot-20261011/tags.txt, source lines 15377--15377. Stable Stacks locator for declared standard input or surrounding context. No selected local proof is replaced by this pointer.\par\medskip
\ifdefined\hypertarget\hypertarget{sourcepointer-2}{}\fi
\noindent\textit{Source reference.} Original lemma 0C8U: focused-seeds/AGP/sources/stacks-snapshot-20261011/tags.txt, source lines 15379--15379. Stable Stacks locator for declared standard input or surrounding context. No selected local proof is replaced by this pointer.\par\medskip
\ifdefined\hypertarget\hypertarget{sourcepointer-3}{}\fi
\noindent\textit{Source reference.} Original lemma 0C8T: focused-seeds/AGP/sources/stacks-snapshot-20261011/tags.txt, source lines 15378--15378. Stable Stacks locator for declared standard input or surrounding context. No selected local proof is replaced by this pointer.\par\medskip
\ifdefined\hypertarget\hypertarget{sourcepointer-4}{}\fi
\noindent\textit{Source reference.} Original lemma 0C5X: focused-seeds/AGP/sources/stacks-snapshot-20261011/tags.txt, source lines 15277--15277. Stable Stacks locator for declared standard input or surrounding context. No selected local proof is replaced by this pointer.\par\medskip
\section{Numerical types}
\label{section-numerical-types}

\noindent
Part of the arguments will involve the combinatorics of the following
data structures.

\begin{definition}
\label{definition-type}
A {\it numerical type} $T$ is given by
$$
n, m_i, a_{ij}, w_i, g_i
$$
where $n \geq 1$ is an integer and $m_i$, $a_{ij}$, $w_i$, $g_i$
are integers for $1 \leq i, j \leq n$ subject to the following conditions
\begin{enumerate}
\item $m_i > 0$, $w_i > 0$, $g_i \geq 0$,
\item the matrix $A = (a_{ij})$ is symmetric and $a_{ij} \geq 0$
for $i \not = j$,
\item there is no proper nonempty subset $I \subset \{1, \ldots, n\}$
such that $a_{ij} = 0$ for $i \in I$, $j \not \in I$,
\item for each $i$ we have $\sum_j a_{ij}m_j = 0$, and
\item $w_i | a_{ij}$.
\end{enumerate}
\end{definition}

\noindent
This is obviously a somewhat annoying type of structure to work with,
but it is exactly what shows up in special fibres of proper regular
models of smooth geometrically connected curves.
Of course we only care about these types up to reordering the indices.

\begin{definition}
\label{definition-type-equivalent}
We say two numerical types $n, m_i, a_{ij}, w_i, g_i$ and
$n', m'_i, a'_{ij}, w'_i, g'_i$ are {\it equivalent types} if
there exists a permutation $\sigma$ of $\{1, \ldots, n\}$
such that $m_i = m'_{\sigma(i)}$, $a_{ij} = a'_{\sigma(i)\sigma(j)}$,
$w_i = w'_{\sigma(i)}$, and $g_i = g'_{\sigma(i)}$.
\end{definition}

\noindent
A numerical type has a genus.

\begin{lemma}
\label{lemma-genus}
Let $n, m_i, a_{ij}, w_i, g_i$ be a numerical type. Then the expression
$$
g = 1 + \sum m_i(w_i(g_i - 1) - \frac{1}{2} a_{ii})
$$
is an integer.
\end{lemma}

\begin{proof}
To prove $g$ is an integer we have to show that $\sum a_{ii}m_i$ is even.
This we can see by computing modulo $2$ as follows
\begin{align*}
\sum\nolimits_i a_{ii} m_i 
& \equiv
\sum\nolimits_{i,\ m_i\text{ odd}} a_{ii}m_i \\
& \equiv
\sum\nolimits_{i,\ m_i\text{ odd}} \sum\nolimits_{j \not = i} a_{ij}m_j \\
& \equiv
\sum\nolimits_{i,\ m_i\text{ odd}}
\sum\nolimits_{j \not = i,\ m_j\text{ odd}} a_{ij}m_j \\
& \equiv
\sum\nolimits_{i < j,\ m_i\text{ and }m_j\text{ odd}} a_{ij}(m_i + m_j) \\
& \equiv
0
\end{align*}
where we have used that $a_{ij} = a_{ji}$ and that $\sum_j a_{ij}m_j = 0$
for all $i$.
\end{proof}

\begin{definition}
\label{definition-genus}
We say $n, m_i, a_{ij}, w_i, g_i$ is a {\it numerical type of genus $g$}
if $g = 1 + \sum m_i(w_i(g_i - 1) - \frac{1}{2} a_{ii})$ is the integer
from Lemma \ref{lemma-genus}.
\end{definition}

\noindent
We will prove below (Lemma \ref{lemma-genus-nonnegative}) that the genus
is almost always $\geq 0$. But you can have numerical types with
negative genus.

\begin{lemma}
\label{lemma-irreducible}
Let $n, m_i, a_{ij}, w_i, g_i$ be a numerical type of genus $g$.
If $n = 1$, then $a_{11} = 0$ and $g = 1 + m_1w_1(g_1 - 1)$.
Moreover, we can classify all such numerical types as follows
\begin{enumerate}
\item If $g < 0$, then $g_1 = 0$ and there are finitely many possible
numerical types of genus $g$ with $n = 1$ corresponding to factorizations
$m_1w_1 = 1 - g$.
\item If $g = 0$, then $m_1 = 1$, $w_1 = 1$, $g_1 = 0$
as in Lemma \ref{lemma-genus-zero}.
\item If $g = 1$, then we conclude $g_1 = 1$ but $m_1, w_1$ can be arbitrary
positive integers; this is case
(\ref{item-one}) of Lemma \ref{lemma-genus-one}.
\item If $g > 1$, then $g_1 > 1$ and there are finitely many possible
numerical types of genus $g$ with $n = 1$ corresponding to
factorizations $m_1w_1(g_1 - 1) = g - 1$.
\end{enumerate}
\end{lemma}

\begin{proof}
The lemma proves itself.
\end{proof}

\begin{lemma}
\label{lemma-diagonal-negative}
Let $n, m_i, a_{ij}, w_i, g_i$ be a numerical type of genus $g$.
If $n > 1$, then $a_{ii} < 0$ for all $i$.
\end{lemma}

\begin{proof}
Lemma \ref{lemma-recurring-symmetric-real} applies to the matrix $A$.
\end{proof}

\begin{lemma}
\label{lemma-minus-one}
Let $n, m_i, a_{ij}, w_i, g_i$ be a numerical type of genus $g$.
Assume $n > 1$. If $i$ is such that the contribution
$m_i(w_i(g_i - 1) - \frac{1}{2} a_{ii})$
to the genus $g$ is $< 0$, then $g_i = 0$ and $a_{ii} = -w_i$.
\end{lemma}

\begin{proof}
Follows immediately from Lemma \ref{lemma-diagonal-negative} and
$w_i > 0$, $g_i \geq 0$, and $w_i | a_{ii}$.
\end{proof}

\begin{definition}
\label{definition-type-minus-one}
Let $n, m_i, a_{ij}, w_i, g_i$ be a numerical type.
We say $i$ is a {\it $(-1)$-index} if $g_i = 0$ and $a_{ii} = -w_i$.
\end{definition}

\noindent
We can ``contract'' $(-1)$-indices.

\begin{lemma}
\label{lemma-contract}
Let $n, m_i, a_{ij}, w_i, g_i$ be a numerical type $T$.
Assume $n$ is a $(-1)$-index. Then there is a numerical
type $T'$ given by $n', m'_i, a'_{ij}, w'_i, g'_i$ with
\begin{enumerate}
\item $n' = n - 1$,
\item $m'_i = m_i$,
\item $a'_{ij} = a_{ij} - a_{in}a_{jn}/a_{nn}$,
\item $w'_i = w_i/2$ if $a_{in}/w_n$ even and $a_{in}/w_i$ odd
and $w'_i = w_i$ else,
\item $g'_i =
\frac{w_i}{w'_i}(g_i - 1) + 1 + \frac{a_{in}^2 - w_na_{in}}{2w'_iw_n}$.
\end{enumerate}
Moreover, we have $g = g'$.
\end{lemma}

\begin{proof}
Observe that $n > 1$ for example by Lemma \ref{lemma-irreducible}
and hence $n' \geq 1$. We check conditions (1) -- (5) of
Definition \ref{definition-type} for $n', m'_i, a'_{ij}, w'_i, g'_i$.

\medskip\noindent
Condition (1) is immediate.

\medskip\noindent
Condition (2). Symmetry of $A' = (a'_{ij})$ is immediate
and since $a_{nn} < 0$ by Lemma \ref{lemma-diagonal-negative}
we see that $a'_{ij} \geq a_{ij} \geq 0$ if $i \not = j$.

\medskip\noindent
Condition (3). Suppose that $I \subset \{1, \ldots, n - 1\}$ such that
$a'_{ii'} = 0$ for $i \in I$ and $i' \in \{1, \ldots, n - 1\} \setminus I$.
Then we see that for each $i \in I$ and $i' \in I'$ we have
$a_{in}a_{i'n} = 0$. Thus either $a_{in} = 0$ for all $i \in I$ and
$I \subset \{1, \ldots, n\}$ is a contradiction for property (3) for $T$,
or $a_{i'n} = 0$ for all $i' \in \{1, \ldots, n - 1\} \setminus I$
and $I \cup \{n\} \subset \{1, \ldots, n\}$ is a contradiction for
property (3) of $T$. Hence (3) holds for $T'$.

\medskip\noindent
Condition (4). We compute
$$
\sum\nolimits_{j = 1}^{n - 1} a'_{ij}m_j =
\sum\nolimits_{j = 1}^{n - 1}
(a_{ij}m_j - \frac{a_{in}a_{jn}m_j}{a_{nn}}) =
- a_{in}m_n - \frac{a_{in}}{a_{nn}}(-a_{nn}m_n) = 0
$$
as desired.

\medskip\noindent
Condition (5). We have to show that $w'_i$ divides $a_{in}a_{jn}/a_{nn}$.
This is clear because $a_{nn} = -w_n$ and $w_n | a_{jn}$ and $w_i | a_{in}$.

\medskip\noindent
To show that $g = g'$ we first write
\begin{align*}
g
& =
1 + \sum\nolimits_{i = 1}^n m_i(w_i(g_i - 1) - \frac{1}{2}a_{ii}) \\
& =
1 + \sum\nolimits_{i = 1}^{n - 1} m_i(w_i(g_i - 1) - \frac{1}{2}a_{ii})
-\frac{1}{2}m_nw_n \\
& =
1 +  \sum\nolimits_{i = 1}^{n - 1}
m_i(w_i(g_i - 1) - \frac{1}{2}a_{ii} - \frac{1}{2}a_{in})
\end{align*}
Comparing with the expression for $g'$ we see that it suffices if
$$
w'_i(g'_i - 1) - \frac{1}{2}a'_{ii} =
w_i(g_i - 1) - \frac{1}{2}a_{in} - \frac{1}{2}a_{ii}
$$
for $i \leq n - 1$. In other words, we have
$$
g'_i = \frac{2w_i(g_i - 1) - a_{in} - a_{ii} + a'_{ii} + 2w'_i}{2w'_i} =
\frac{w_i}{w'_i}(g_i - 1) + 1 + \frac{a_{in}^2 - w_na_{in}}{2w'_iw_n}
$$
It is elementary to check that this is an integer $\geq 0$
if we choose $w'_i$ as in (4).
\end{proof}

\begin{lemma}
\label{lemma-top-genus}
Let $n, m_i, a_{ij}, w_i, g_i$ be a numerical type.
Let $e$ be the number of pairs $(i, j)$ with $i < j$ and $a_{ij} > 0$.
Then the expression $g_{top} = 1 - n + e$ is $\geq 0$.
\end{lemma}

\begin{proof}
If not, then $e < n - 1$ which means there exists an $i$ such that
$a_{ij} = 0$ for all $j \not = i$. This contradicts assumption
(3) of Definition \ref{definition-type}.
\end{proof}

\begin{definition}
\label{definition-top-genus}
Let $n, m_i, a_{ij}, w_i, g_i$ be a numerical type $T$. The
{\it topological genus of $T$} is the nonnegative integer
$g_{top} = 1 - n + e$ from Lemma \ref{lemma-top-genus}.
\end{definition}

\noindent
We want to bound the genus by the topological genus. However, this
will not always be the case, for example for numerical types
with $n = 1$ as in Lemma \ref{lemma-irreducible}. But it will
be true for minimal numerical types which are defined as follows.

\begin{definition}
\label{definition-type-minimal}
We say the numerical type $n, m_i, a_{ij}, w_i, g_i$ of genus $g$
is {\it minimal} if there does not exist an $i$
with $g_i = 0$ and $a_{ii} = -w_i$, in other words, if there
does not exist a $(-1)$-index.
\end{definition}

\noindent
We will prove that the genus $g$ of a minimal type with $n > 1$
is greater than or equal to $\max(1, g_{top})$.


\begin{lemma}
\label{lemma-non-irreducible-minimal-type-genus-at-least-one}
If $n, m_i, a_{ij}, w_i, g_i$ is a minimal numerical type
with $n > 1$, then $g \geq 1$.
\end{lemma}

\begin{proof}
This is true because $g = 1 + \sum \Phi_i$ with
$\Phi_i = m_i(w_i(g_i - 1) - \frac{1}{2} a_{ii})$ nonnegative
by Lemma \ref{lemma-minus-one} and the definition of minimal types.
\end{proof}

\begin{lemma}
\label{lemma-genus-nonnegative}
If $n, m_i, a_{ij}, w_i, g_i$ is a minimal numerical type
with $n > 1$, then $g \geq g_{top}$.
\end{lemma}

\begin{proof}
The reader who is only interested in the case of numerical types
associated to proper regular models can skip this proof as we will
reprove this in the geometric situation later.
We can write
$$
g_{top} = 1 - n + \frac{1}{2}\sum\nolimits_{a_{ij} > 0} 1 =
1 + \sum\nolimits_i (-1 +
\frac{1}{2}\sum\nolimits_{j \not = i,\ a_{ij} > 0} 1) 
$$
On the other hand, we have
\begin{align*}
g & =
1 + \sum m_i(w_i(g_i - 1) - \frac{1}{2} a_{ii}) \\
& =
1 + \sum m_iw_ig_i - \sum m_iw_i +
\frac{1}{2} \sum\nolimits_{i \not = j} a_{ij}m_j \\
& =
1 + \sum\nolimits_i
m_iw_i(-1 + g_i + \frac{1}{2} \sum\nolimits_{j \not = i} \frac{a_{ij}}{w_i})
\end{align*}
The first equality is the definition, the second equality uses that
$\sum a_{ij}m_j = 0$, and the last equality uses that
uses $a_{ij} = a_{ji}$ and switching order of
summation. Comparing with the formula for $g_{top}$ we conclude
that the lemma holds if
$$
\Psi_i =
m_iw_i(-1 + g_i + \frac{1}{2} \sum\nolimits_{j \not = i} \frac{a_{ij}}{w_i})
- (-1 + \frac{1}{2}\sum\nolimits_{j \not = i,\ a_{ij} > 0} 1)
$$
is $\geq 0$ for each $i$. However, this may not be the case.
Let us analyze for which indices we can have $\Psi_i < 0$.
First, observe that
$$
(-1 + g_i + \frac{1}{2}\sum\nolimits_{j \not = i} \frac{a_{ij}}{w_i}) \geq
(-1 + \frac{1}{2}\sum\nolimits_{j \not = i,\ a_{ij} > 0} 1)
$$
because $a_{ij}/w_i$ is a nonnegative integer. Since $m_iw_i$ is
a positive integer we conclude that $\Psi_i \geq 0$ as soon as
either $m_iw_i = 1$ or the left hand side of the inequality is $\geq 0$
which happens if $g_i > 0$, or $a_{ij} > 0$ for at least two indices $j$, or
if there is a $j$ with $a_{ij} > w_i$. Thus
$$
P = \{i : \Psi_i < 0\}
$$
is the set of indices $i$ such that $m_iw_i > 1$, $g_i = 0$,
$a_{ij} > 0$ for a unique $j$, and $a_{ij} = w_i$ for this $j$.
Moreover
$$
i \in P \Rightarrow \Psi_i = \frac{1}{2}(-m_iw_i + 1)
$$
The strategy of proof is to show that given $i \in P$ we can borrow a bit
from $\Psi_j$ where $j$ is the neighbour of $i$, i.e., $a_{ij} > 0$.
However, this won't quite work because $j$ may be an index with $\Psi_j = 0$.

\medskip\noindent
Consider the set
$$
Z = \{j : g_j = 0\text{ and }
j\text{ has exactly two neighbours }i, k\text{ with }
a_{ij} = w_j = a_{jk}\}
$$
For $j \in Z$ we have $\Psi_j = 0$. We will consider sequences
$M = (i, j_1, \ldots, j_s)$ where $s \geq 0$,
$i \in P$, $j_1, \ldots, j_s \in Z$, and
$a_{ij_1} > 0, a_{j_1j_2} > 0, \ldots, a_{j_{s - 1}j_s} > 0$.
If our numerical type consists of two indices which are in $P$
or more generally if our numerical type consists of
two indices which are in $P$ and all other indices in $Z$, then
$g_{top} = 0$ and we win by
Lemma \ref{lemma-non-irreducible-minimal-type-genus-at-least-one}.
We may and do discard these cases.

\medskip\noindent
Let $M = (i, j_1, \ldots, j_s)$ be a maximal sequence and let
$k$ be the second neighbour of $j_s$. (If $s = 0$, then $k$
is the unique neighbour of $i$.) By maximality $k \not \in Z$
and by what we just said $k \not \in P$. Observe that
$w_i = a_{ij_1} = w_{j_1} = a_{j_1j_2} = \ldots = w_{j_s} =
a_{j_sk}$. Looking at the definition
of a numerical type we see that
\begin{align*}
m_ia_{ii} + m_{j_1}w_i & = 0,\\
m_iw_i + m_{j_1}a_{j_1j_1} + m_{j_2}w_i & = 0,\\
\ldots & \ldots \\
m_{j_{s - 1}}w_i + m_{j_s}a_{j_sj_s} + m_kw_i & = 0
\end{align*}
The first equality implies $m_{j_1} \geq 2m_i$ because the
numerical type is minimal. Then the second equality implies
$m_{j_2} \geq 3m_i$, and so on. In any case, we conclude that
$m_k \geq 2m_i$ (including when $s = 0$).

\medskip\noindent
Let $k$ be an index such that we have a $t > 0$ and pairwise distinct
maximal sequences $M_1, \ldots, M_t$ as above, with
$M_b = (i_b, j_{b, 1}, \ldots, j_{b, s_b})$
such that $k$ is a neighbour of $j_{b, s_b}$ for $b = 1, \ldots, t$.
We will show that $\Phi_j + \sum_{b = 1, \ldots, t} \Phi_{i_b} \geq 0$.
This will finish the proof of the lemma by what we said above.
Let $M$ be the union of the indices occurring in $M_b$, $b = 1, \ldots, t$.
We write
$$
\Psi_k =
-\sum\nolimits_{b = 1, \ldots, t} \Psi_{i_b} + \Psi_k'
$$
where
\begin{align*}
\Psi_k' & =
m_kw_k\left(-1 + g_k +
\frac{1}{2} \sum\nolimits_{b = 1, \ldots t}
(\frac{a_{kj_{b, s_b}}}{w_k} - \frac{m_{i_b}w_{i_b}}{m_kw_k}) +
\frac{1}{2} \sum\nolimits_{l \not = k,\ l \not \in M}
\frac{a_{kl}}{w_k}
\right) \\
&
-\left(
-1 + \frac{1}{2}\sum\nolimits_{l \not = k,\ l \not \in M,\ a_{kl} > 0} 1
\right)
\end{align*}
Assume $\Psi_k' < 0$ to get a contradiction.
If the set $\{l : l \not = k,\ l \not \in M,\ a_{kl} > 0\}$ is empty,
then $\{1, \ldots, n\} = M \cup \{k\}$ and $g_{top} = 0$
because $e = n - 1$ in this case and the result holds by
Lemma \ref{lemma-non-irreducible-minimal-type-genus-at-least-one}.
Thus we may assume there is at least one such $l$ which contributes
$(1/2)a_{kl}/w_k \geq 1/2$ to the sum inside the first brackets.
For each $b = 1, \ldots, t$ we have
$$
\frac{a_{kj_{b, s_b}}}{w_k} - \frac{m_{i_b}w_{i_b}}{m_kw_k} =
\frac{w_{i_b}}{w_k}(1 - \frac{m_{i_b}}{m_k})
$$
This expression is $\geq \frac{1}{2}$ because $m_k \geq 2m_{i_b}$
by the previous paragraph and is $\geq 1$ if $w_k < w_{i_b}$.
It follows that $\Psi_k' < 0$ implies $g_k = 0$.
If $t \geq 2$ or $t = 1$ and $w_k < w_{i_1}$, then $\Psi_k' \geq 0$
(here we use the existence of an $l$ as shown above) which
is a contradiction too.
Thus $t = 1$ and $w_k = w_{i_1}$. If there at least two nonzero terms
in the sum over $l$ or if there is one such $k$ and $a_{kl} > w_k$, then
$\Psi_k' \geq 0$ as well. The final possibility is that $t = 1$ and
there is one $l$ with $a_{kl} = w_k$. This is disallowed as this would
mean $k \in Z$ contradicting the maximality of $M_1$.
\end{proof}

\bibliographystyle{plain}
\bibliography{my}
\end{document}
